\subsection{Local improvement without a completed root reference}
\label{sec:frontier}
The root comparison can exhaust its budget before valuing the parent, even
when later decisions admit tractable local comparisons. We address this
failure by following the parent to a public decision boundary and replacing
only continuations with a certified local gain. The original controller's
completed selections are preserved. Figure~\ref{fig:frontier-overview} shows
the shared prefix and the continuations that can change.

The boundary is the learner's next discard. We first enumerate the entire
parent continuation to that decision or to termination, whichever comes first.
All worlds that arrive at the same own/public history form one conditional
belief. Completing this prefix is essential: an unfinished branch may still
reach an already encountered history and change its posterior. An incomplete
prefix therefore supplies no usable local belief.

For each complete frontier belief, search proposes a complete continuation and
compares its lower value bound with the unchanged parent's upper bound.
Unresolved probability mass contributes the interval $[0,1]$ for terminal
win score. A strict separation certifies improvement even if the local parent
value remains unresolved. The program uses explicit legal overrides and a
parent default. On reaching the frontier, execution commits to that one local
program or keeps the parent; it never switches between sibling programs.

\begin{figure}[htbp]
\centering
\begingroup
\footnotesize
\setlength{\fboxsep}{5pt}
\fbox{\begin{minipage}{0.89\linewidth}
\centering
\textbf{Complete the common parent prefix}\\[2pt]
Enumerate every compatible world to the learner's next discard or termination.\\
Group arrivals by own/public history; retain each group's full belief and mass.
\end{minipage}}

\smallskip
\begin{tabular}{@{}c@{\hspace{0.025\linewidth}}c@{\hspace{0.025\linewidth}}c@{}}
$\downarrow$ certified histories & $\downarrow$ other histories & $\downarrow$ terminal paths \\[3pt]
\fbox{\begin{minipage}[c][8.5em][c]{0.27\linewidth}
\centering
\textbf{Adopt one local program}\\[3pt]
Candidate lower bound exceeds parent upper bound.\\[3pt]
Strictly positive contribution.
\end{minipage}}
&
\fbox{\begin{minipage}[c][8.5em][c]{0.27\linewidth}
\centering
\textbf{Keep the parent}\\[3pt]
No certified separation, or no budget for this history.\\[3pt]
Gain $=0$ by identical continuation.
\end{minipage}}
&
\fbox{\begin{minipage}[c][8.5em][c]{0.27\linewidth}
\centering
\textbf{Keep the outcome}\\[3pt]
The round ended before the next learner discard.\\[3pt]
Gain $=0$ by identical outcome.
\end{minipage}}
\end{tabular}
\endgroup
\caption{Frontier extension after an unresolved parent reference. The common
prefix partitions all probability mass before any local comparison. Certified
gains add across disjoint arrival histories; unchanged continuations need no
value calculation. Execution follows the common prefix and commits to the
single program for the reached history. If the prefix is incomplete, the
extension keeps the parent everywhere.}
\label{fig:frontier-overview}
\end{figure}


\paragraph{Conditional gain.}
Fix a root belief of integer mass $M$, the opponent model and the chance law.
Let $\mathcal A$ contain the adopted frontier histories, $M_h$ their arrival
masses, and $L_h>U_h$ the candidate-lower and parent-upper comparisons. Then
\begin{equation}
 V(\pi_{\mathrm{patch}})-V(\pi)
 \;\geq\;
 \sum_{h\in\mathcal A}\frac{M_h}{M}(L_h-U_h)\;\geq\;0.
 \label{eq:frontier-gain}
\end{equation}
Both policies follow the same prefix, giving identical arrival probabilities
and complete conditional beliefs. First-arrival events are disjoint. Paths
that terminate earlier or reach an unpatched history have identical
continuations. Conditioning on these events proves the inequality. Full
independent values for the adopted local programs give the exact increment
by the same decomposition, without valuing unchanged paths.

The extension uses the original controller's otherwise unused second-phase
allowance after a failed parent reference, where the original controller
already falls back to the parent. One shared two-million-transition
cap covers prefix construction and every local attempt. At the five-draw root,
the next-discard prefix needs at most 16,380 rows and 131,040 transitions,
within the existing prefix limits. Roots with fewer draws retain the original
controller. Appendix~\ref{app:frontier} gives the bound, scheduling rules and
execution checks. The maximum allowance is unchanged; actual work increases
when the additional phase runs.

In a fixed development screen, the extension recovers improvements at
\FrontierPositiveRoots{} of \FrontierRootCount{} roots whose parent reference
could not be completed. It adopts \FrontierAdoptedNodes{} local programs.
Independent Python replay verifies their conditional gains and every complete
frontier belief. The retained actual trajectories are unchanged; improvements
occur on other supported continuations. These cases demonstrate recovery of
usable improvements. The 8,192-source population result above evaluates the
base controller, without this extension.

\paragraph{Incremental population result.}
The separate \FrontierIncrementSources{}-source study compares the extension
with the base controller. Sources follow the same independent uniform-deal
and seat law as Section~\ref{sec:uniform-results}, with 1,024 matches in each
parent--opponent population. Every source remains in the denominator. The
comparison uses the exact increment from independent terminal replay at
adopted histories, and zero where the complete policies coincide.
The fixed primary criterion requires positive simultaneous lower bounds for
both equal-opponent parent mixtures. Table~\ref{tab:frontier-increment} reports
the complete cohort; Appendix~\ref{app:frontier} records the scope correction,
all source outcomes and computational work.

\begin{table}[htbp]
\centering\small
% Complete incremental study: hybrid minus existing controller.
\begin{tabular}{llrrr}
\toprule
Parent & Opponent & Sources & Mean (pp) & Lower bound (pp) \\
\midrule
1 & WP & 1,024 & $6.77034\times 10^{-4}$ & $6.64356\times 10^{-5}$ \\
1 & Point & 1,024 & 0.00189203 & $2.73724\times 10^{-5}$ \\
2 & WP & 1,024 & 0.00106778 & $1.10018\times 10^{-5}$ \\
2 & Point & 1,024 & $5.00209\times 10^{-4}$ & $7.94267\times 10^{-6}$ \\
\midrule
1 & Equal mixture & 2,048 & 0.00128453 & $4.69040\times 10^{-5}$ \\
2 & Equal mixture & 2,048 & $7.83993\times 10^{-4}$ & $9.47225\times 10^{-6}$ \\
\bottomrule
\end{tabular}

\caption{Frontier extension minus base controller: expected terminal win-score
increments in percentage points. Bounds are simultaneous at 95\% across the
four populations; the fixed equal-weight mixtures inherit that coverage.
All 4,096 sources contribute, including missing roots and unchanged policies.
Lower bounds are rounded downward.}
\label{tab:frontier-increment}
\end{table}

The primary criterion is \FrontierIncrementCriterion{}. All four individual
population lower bounds are also positive. The extension improves 26 source
beliefs, with equal-opponent means of \FrontierIncrementMeanOne{} and
\FrontierIncrementMeanTwo{} percentage points. Thus a failed root reference
can conceal a recoverable improvement in naturally reached play. The measured
increment is small, and realized match-score differences are zero in every
population. The earlier base-controller 0.1-point criterion remains a separate
failed test. Additional planning uses 111,874,298 physical transitions across
this cohort, within the unchanged maximum per-root allowance.
